\begin{table}[H]
\scriptsize
\setlength{\tabcolsep}{3pt}
\renewcommand{\arraystretch}{0.95}

\caption{\label{tab:ea20-counterfactual-income-gap-product-contributions}Household price policies: product-group contributions to no-policy inflation inequality in the 20-country euro area, June 2021--June 2023}
\centering
\resizebox{\textwidth}{!}{%
\begin{tabular}{lrrrrr}
\toprule
Product group & \shortstack{Counterfactual\\ inflation} & \shortstack{Q1 counterfactual\\ contribution} & \shortstack{Q5 counterfactual\\ contribution} & \shortstack{Counterfactual\\ inflation gap} & \shortstack{Effect of price policies\\ on inflation inequality}\\
\midrule
Food and non-alcoholic beverages & 4.6 & 5.0 & 4.4 & 0.7 & 0.1\\
Alcoholic beverage, tobacco and narcotics & 0.5 & 0.6 & 0.4 & 0.2 & 0.0\\
Clothing and footwear & 0.3 & 0.3 & 0.4 & -0.2 & 0.0\\
Housing (rentals and repairs) & 0.5 & 0.4 & 0.6 & -0.2 & 0.0\\
Water, electricity, gas and other fuels & 7.4 & 8.8 & 6.2 & 2.6 & -1.0\\
\addlinespace
Transport & 2.3 & 1.7 & 3.0 & -1.3 & 0.2\\
Other & 0.3 & 0.2 & 0.4 & -0.2 & 0.0\\
\midrule
\textbf{Total} & \textbf{16.0} & \textbf{17.0} & \textbf{15.4} & \textbf{1.6} & \textbf{-0.7}\\
\bottomrule
\end{tabular}%
}
\vspace{-0.4em}
\begin{flushleft}
\footnotesize{Notes. The no-policy counterfactual replaces observed prices for affected COICOP components with estimated paths without household price policies and otherwise retains observed prices. Product-group contributions use RAS expenditure weights at COICOP level 2, excluding imputed rent. The euro-area aggregate covers all 20 countries and uses annual HICP country weights. The inflation gap is Q1 minus Q5. The policy effect is the observed product-group gap minus its no-policy value; a negative number therefore indicates that household price policies reduced the gap. Income transfers and firm-only measures are excluded. All entries are percentage-point contributions. Sources: Eurostat, national statistical institutes, authors' calculations.}
\end{flushleft}
\end{table}
